\documentclass[../main.tex]{subfiles}
\begin{document}
\section{Seguridad en la Empresa}

Esta unidad cambia el foco: deja la matemática y las vulnerabilidades de
aplicación y mira cómo una \textbf{organización real} (corporación o PyME
avanzada) se protege. El profe aclaró que \textbf{este tema no está en Bishop}
(salvo lo que se pueda leer suelto); la clase está muy enfocada en el panorama
profesional contemporáneo y en el estándar \textbf{NIST 800-207} (Zero Trust).

\prof{Esto de clase que vamos a ver hoy no está en Bishop. Está más enfocado en el panorama contemporáneo, lo profesional, cómo se maneja esto en una organización.}

\begin{examen}
La lectura recomendada de esta clase es \textbf{NIST 800-207 (capítulos 2 y 3)},
no Bishop. Toda la parte de Zero Trust sigue ese estándar.
\end{examen}

% =====================================================================
\subsection{Gobernanza (Governance)}

\begin{definicion}{Gobernanza de seguridad}
Conjunto de \textbf{procesos, políticas, procedimientos y controles} que se
implementan para gestionar y mitigar los riesgos de la organización, y el
establecimiento de \textbf{responsabilidades y roles} para garantizar la
protección de la información. Es como el ``gobierno'' interno: emite las reglas
y se hace responsable de que se cumplan.
\end{definicion}

La gobernanza la lleva (generalmente) el equipo de seguridad. Su propósito es
asegurar \textbf{confidencialidad, integridad y disponibilidad} de los datos,
tanto de forma \textbf{proactiva} (poner barreras y controles) como
\textbf{reactiva} (qué medidas tomo una vez atacado).

\textbf{Objetivos principales:}
\begin{itemize}
  \item \textbf{Alinear la estrategia de seguridad con los objetivos del
    negocio.} No poner medidas que vayan en contra de cómo la empresa gana
    dinero.
  \item Reducir y mitigar los riesgos de la operación.
  \item Cumplir regulaciones y normativas aplicables (p.ej.\ empresas que
    cotizan en bolsa deben informar incidentes; datos personales mal protegidos
    $\Rightarrow$ multas millonarias).
  \item \textbf{Crear cultura de seguridad} (campañas de concientización, de
    phishing, capacitación), porque el eslabón más débil es el usuario.
\end{itemize}

\begin{destacado}
El eslabón más débil es siempre el \textbf{usuario / empleado / cliente}. Por
eso, si se ponen \textbf{demasiadas trabas}, el empleado busca atajos (se baja
los datos a un CSV, los sube a Drive, etc.) y \textbf{termina atentando contra
la seguridad}. Esto es el problema de la \emph{aceptación psicológica}: la
estrategia debe estar alineada con el negocio para no friccionar el trabajo.
\end{destacado}

% =====================================================================
\subsection{Gestión de riesgos (Risk Management)}

Históricamente la seguridad se gestionaba ``como un riesgo más'', y ese modelo
se mantiene. Es un \textbf{ciclo que nunca para} (los riesgos evolucionan,
aparecen nuevos): \textbf{Identificar $\to$ Evaluar (Assess) $\to$ Mitigar
$\to$ Monitorear} y vuelta a empezar.

\subsubsection{Identificación del riesgo}
Primero se identifican los \textbf{activos críticos}, que involucran:
\textbf{Datos} (de empleados, clientes, comerciales, financieros, proveedores),
\textbf{Sistemas}, \textbf{Infraestructura} (servidores, oficinas, data
centers) y \textbf{Personas}.

Luego los \textbf{riesgos potenciales}: Malware, Ataques externos, Errores
humanos, Sabotaje (interno) y Ambientales (huracán, incendio). El profe sumó
los \textbf{problemas externos} (caída de una región de un proveedor cloud que
voltea a muchas empresas a la vez).

\subsubsection{Evaluación del riesgo}
Se estima \textbf{probabilidad de ocurrencia} $\times$ \textbf{impacto en el
negocio} (matriz probabilidad vs.\ impacto). Se clasifican en \textbf{Alto /
Medio / Bajo} y se atacan primero los de mayor probabilidad y mayor impacto.
Acá también se realizan evaluaciones externas (Red Team / PenTest).

\subsubsection{Mitigación del riesgo --- las 5 opciones}
\begin{definicion}{Tratamiento de un riesgo}
Ante un riesgo identificado hay cinco maneras de manejarlo:
\begin{enumerate}
  \item \textbf{Evitarlo}: eliminar la causa. Suele ser \emph{costoso/complejo}
    (p.ej.\ resolver vulnerabilidades de un sistema legacy o migrarlo).
  \item \textbf{Reducirlo} (lo más común): bajar el impacto sin eliminarlo
    (p.ej.\ sumar un factor extra de autenticación a una autenticación débil).
  \item \textbf{Transferirlo}: delegarlo a un proveedor (p.ej.\ Cloudflare como
    proxy reverso). Si pasa algo, responde el proveedor.
  \item \textbf{Compensarlo}: no se puede reducir, pero se contrarresta el
    impacto (p.ej.\ auditar/controlar continuamente una dependencia que no se
    puede quitar).
  \item \textbf{Aceptarlo}: se entiende y se asume el riesgo (típico en riesgos
    de baja probabilidad). \textbf{Se debe documentar} la decisión.
\end{enumerate}
\end{definicion}

\begin{examen}
Distinguir \textbf{reducir} (baja la probabilidad/impacto del riesgo) de
\textbf{compensar} (el riesgo sigue igual, pero se contrarresta su efecto). Un
ejemplo de \textbf{aceptar}: muchas empresas tienen un \textbf{DRP} (Disaster
Recovery Plan) documentado y exigido (p.ej.\ el Banco Central a las fintech),
pero \textbf{no lo ejecutan} por costo $\Rightarrow$ aceptan el riesgo y lo
dejan documentado.
\end{examen}

\subsubsection{Monitoreo}
Es cíclico y permanente (la organización es un sistema vivo). Existen
\textbf{frameworks} de apoyo: NIST (análogo del IRAM en Argentina) y los
\textbf{CIS Controls} (18 controles) como guía de qué tener en cuenta.

% =====================================================================
\subsection{Políticas de seguridad}

\begin{definicion}{Política de seguridad}
Documento \textbf{genérico y muy abarcativo} (``la Biblia'' de la seguridad de
la empresa); se confecciona cada cierto tiempo y lo aprueba la dirección
(CISO/CEO). Define todo lo que compete a la seguridad, llegando incluso al plano
legal (manejo de datos personales, estándares de industria).
\end{definicion}

De la política se va ``bajando línea'' con creciente especificidad (y decreciente
\emph{ownership} del management). Pirámide:
\begin{itemize}
  \item \textbf{Policy} (declaración general del management) --- mandatorio.
  \item \textbf{Standards} (controles específicos mandatorios).
  \item \textbf{Procedures} (instrucciones paso a paso).
  \item \textbf{Guidelines} (recomendaciones / mejores prácticas).
  \item \textbf{Baselines} (formas uniformes de implementación; cursos,
    capacitaciones, concientización).
\end{itemize}

\begin{ojo}
En la pirámide, hacia \textbf{arriba} crece el \textbf{ownership del management}
y es \textbf{mandatorio}; hacia \textbf{abajo} crece la \textbf{especificidad}
(y se vuelve más discrecional). No confundir el eje.
\end{ojo}

% =====================================================================
\subsection{Estrategias: Defense-in-depth vs.\ Zero Trust}

Surgieron estrategias genéricas para implementar las políticas. La clase destaca
dos: \textbf{Defense-in-depth} (defensa en profundidad) y \textbf{Zero Trust
Architecture (ZTA)}.

\subsubsection{Defense-in-depth (defensa en profundidad)}
Asume que \textbf{ningún control es perfecto}. Diseña \textbf{capas}, cada una
pensada para contener a la siguiente (modelo ``cebolla''), dando
\textbf{redundancia}: si falla una capa, la de abajo contiene.

Las \textbf{5 capas} (de afuera hacia el dato):
\begin{enumerate}
  \item \textbf{Perímetro (Perimeter)}: frontera con Internet. Firewall que
    deniega lo que viene de afuera y solo acepta cierta red o VPNs; honeypot;
    seguridad de mensajería (anti-virus/anti-malware); IDS/IPS de perímetro.
  \item \textbf{Red (Network)}: dentro de la red corporativa.
    \textbf{Segmentación} (separar la red de operaciones de la red de
    servidores); NAC; Web Proxy / Content Filtering; DLP.
  \item \textbf{Host / Endpoint}: \emph{hardening} (``jardinizar'') del host
    --- dejarlo inmutable, solo con el software permitido (idea tipo
    contenedor); Host IDS/IPS; antivirus.
  \item \textbf{Aplicación}: buen diseño y buenas prácticas para no meter
    vulnerabilidades; \textbf{WAF}; monitoreo/escaneo de BD.
  \item \textbf{Datos}: lo que más importa. Clasificación, cifrado, gestión de
    identidades y accesos (IAM), DLP, PKI.
\end{enumerate}

Para llegar al dato, el atacante debe pasar perímetro $\to$ moverse
lateralmente a la red de servidores (que no debería ser alcanzable) $\to$
vulnerar el host jardinizado $\to$ vulnerar la aplicación.

\paragraph{Principios.}
\begin{itemize}
  \item \textbf{Diversificación}: distintas tecnologías/métodos $\Rightarrow$ una
    sola vulnerabilidad no atraviesa todo.
  \item \textbf{Redundancia}: cada capa contiene a la siguiente.
  \item \textbf{Segmentación}: dividir la red en segmentos controlados (se puede
    llegar al extremo de que una aplicación solo ``vea'' a la aplicación con la
    que necesita hablar).
  \item \textbf{Monitoreo}: visibilidad de todas las capas y
    \textbf{correlación de eventos} (si no se monitorea, nada de lo anterior
    importa).
\end{itemize}

\paragraph{Problemas.}
\begin{itemize}
  \item \textbf{Altos costos}: de tecnología (muchos proveedores), de
    capacitación/operación y de monitoreo (se recolectan muchos datos pero se
    les saca poco valor).
  \item \textbf{No resuelve bien los ataques internos} (si alguien ya está
    adentro, es difícil detectarlo; muchas herramientas no son bidireccionales).
  \item \textbf{Difícil en infraestructura distribuida}: pensado para
    \emph{on-premise}/data center propio; hoy hay data centers híbridos,
    multi-nube y empresas \emph{cloud native}.
\end{itemize}

\begin{ojo}
\textbf{Firewalls} en esta unidad (según lo visto en clase): el firewall de
\textbf{perímetro} bloquea/filtra tráfico de red; el \textbf{WAF (Web
Application Firewall)} aparece en la capa de \textbf{aplicación} y es ``un
firewall pero para el tráfico de una aplicación web''. En clase se mencionaron
solo estos dos (firewall de perímetro y WAF).
\end{ojo}

\begin{destacado}
\textbf{DMZ (zona desmilitarizada).} En el diagrama de Defense-in-depth
aparecen las \emph{Secure DMZs} dentro de la seguridad de
\textbf{perímetro/red}: es la zona expuesta hacia la red externa.
\end{destacado}

\paragraph{IDS / IPS.}
\begin{definicion}{IDS / IPS}
Sistemas de \textbf{detección} (IDS, \emph{Intrusion Detection System}) y de
\textbf{prevención} (IPS, \emph{Intrusion Prevention System}) de intrusiones.
Funcionan ``tipo antivirus'': detectan patrones a través de \textbf{firmas} o
\textbf{anomalías} y determinan si \textbf{bloquear} (prevenir), \textbf{alertar}
o tomar alguna acción.
\end{definicion}

\begin{examen}
La diferencia clave: \textbf{IDS detecta} (alerta), \textbf{IPS previene}
(bloquea). En el diagrama de Defense-in-depth aparecen tanto en el
\textbf{perímetro} como en el \textbf{host} (Host IDS/IPS).
\end{examen}

\paragraph{Otros componentes vistos en las capas.}
\begin{itemize}
  \item \textbf{Honeypot}: red ``boba'' a donde se deriva al atacante detectado
    en el perímetro para mantenerlo entretenido y que no llegue a la red que
    importa.
  \item \textbf{NAC (Network Access Control)}: define si un dispositivo se puede
    conectar o no (por IP, MAC, firmware, credenciales); similar a limitar por
    MAC en un router.
  \item \textbf{Web Proxy / Content Filtering} (p.ej.\ Fortinet): bloquea/filtra
    tráfico (que la gente no entre a sitios maliciosos o no laborales).
  \item \textbf{DLP (Data Loss Prevention)}: evita la \textbf{exfiltración} de
    datos hacia afuera (escanea por datos de tarjeta, personales, etc.).
\end{itemize}

% =====================================================================
\subsubsection{Zero Trust Architecture (ZTA)}

\begin{definicion}{Zero Trust Architecture (NIST 800-207)}
Modelo de seguridad que asume que las amenazas \textbf{internas y externas están
presentes en la red en todo momento}. Objetivo: proteger integridad,
confidencialidad y disponibilidad de datos y recursos. Principios:
\textbf{Nunca confíes (Never Trust)}, \textbf{Siempre verifica (Always Verify)},
\textbf{Asume que has sido atacado (Assume Breach)}.
\end{definicion}

ZTA nace de constatar que \textbf{ninguna red es segura}: la workstation de un
empleado (en la oficina o por VPN) puede tener malware sin que lo sepa, o puede
haber un intruso físico. Por eso \textbf{cualquier red interna pasa a
considerarse tan peligrosa como Internet}, y los controles que se aplicaban
``afuera'' ahora se aplican también ``adentro''.

\begin{destacado}
La \textbf{intranet ya no es de confianza}. Antes era ``buena práctica'' relajar
controles dentro de la red corporativa (para no friccionar al empleado); ZTA
revierte eso: \textbf{cada componente/aplicación debe forzar la verificación}.
En la práctica se relaja con \textbf{federación de identidad / SSO} (el empleado
loguea una vez, p.ej.\ contra el directorio, y cada aplicación hereda su
identidad, roles y permisos sin pedir login nuevo, pero \emph{sí} aplicando los
controles).
\end{destacado}

\paragraph{Componentes (NIST 800-207).} El núcleo es el \textbf{Policy Decision
Point (PDP)} = \textbf{Policy Engine} + \textbf{Policy Administrator}, y el
\textbf{Policy Enforcement Point (PEP)} que separa la zona \emph{Untrusted}
(Subject/System) de la \emph{Trusted} (Enterprise Resource). Lo alimentan:
\begin{itemize}
  \item \textbf{CDM (Continuous Diagnostics and Mitigation)}: informa al motor
    de políticas sobre el activo que pide acceso (¿antivirus actualizado?,
    ¿últimos updates de seguridad?); puede desactivar el recurso si no cumple.
  \item \textbf{Industry Compliance}: garantiza el cumplimiento de regímenes
    normativos --- \textbf{HIPAA, Banco Central, SOX, PCI}.
  \item \textbf{Threat Intelligence}: fuentes internas/externas sobre CVEs,
    firmas de ataques, IPs comprometidas, botnets (p.ej.\ VirusTotal).
  \item \textbf{Activity Logs}: registros de actividad de red y sistemas (datos
    crudos, en tiempo casi real).
  \item \textbf{SIEM (Security Information and Event Management)}: recolecta y
    \textbf{correlaciona} eventos para detectar anomalías y generar alertas.
  \item \textbf{ID Management}: crea/almacena/gestiona identidades de los
    usuarios (una persona puede tener varios usuarios), con sus roles y
    permisos.
  \item \textbf{PKI}: genera y registra los certificados de sujetos, recursos,
    servicios y aplicaciones.
  \item \textbf{Data Access Policy}: reglas/políticas de acceso a los recursos;
    punto de partida para autorizar (codificadas o generadas dinámicamente por
    el PE).
\end{itemize}

\begin{ojo}
ZTA ``parece la panacea'' pero es \textbf{muy fácil de implementar mal}: o se
implementa de forma deficiente, o se implementa y \textbf{no se controla / no se
hace evolucionar}. Tener ZTA ``puesto'' no protege si no se monitorea y audita.
Si una empresa arranca \textbf{de cero}, conviene ir directo a ZTA (sin el
sesgo/lastre del esquema viejo).
\end{ojo}

% =====================================================================
\subsection{Privileged Access Management (PAM)}

\begin{definicion}{PAM}
Práctica para \textbf{proteger y gestionar cuentas privilegiadas} (acceso a
sistemas críticos). Estos accesos son sensibles porque pueden conceder
capacidades amplias, modificar configuraciones de seguridad, acceder a datos
confidenciales y supervisar otras cuentas. Son lo que más buscan los atacantes.
\end{definicion}

\textbf{Prácticas/herramientas:}
\begin{enumerate}
  \item \textbf{Inventario} de cuentas privilegiadas (visibilidad y control).
  \item \textbf{Privilegio mínimo}: la cuenta privilegiada \textbf{no} es la
    cuenta de trabajo de la persona; es una cuenta separada (que puede
    transferirse al cambiar de rol).
  \item \textbf{MFA} (autenticación multifactor) para garantizar acceso solo de
    autorizados y el \textbf{no repudio}.
  \item \textbf{Registro y monitoreo} de toda la actividad (logs en servidor
    externo para reconstruir qué pasó).
  \item \textbf{Rotación de contraseñas} (regular y automática; y reválida al
    transferir la cuenta).
\end{enumerate}
Ciclo de vida PAM: Define $\to$ Discover $\to$ Manage \& Protect $\to$ Monitor
$\to$ Detect Abnormal Usage $\to$ Respond to Incidents $\to$ Review \& Audit.

% =====================================================================
\subsection{Prevención de vulnerabilidades: DevSecOps}

\begin{definicion}{DevSecOps}
Metodología que mete la \textbf{seguridad en cada paso} del proceso de llevar
algo a producción, integrándola dentro de los equipos de desarrollo. Es la
evolución de \textbf{DevOps} (que unió Dev y Ops bajo un mismo objetivo:
``\emph{you build it, you run it}'') sumándole el equipo de seguridad.
\end{definicion}

El problema histórico: Dev quiere sacar features rápido; Seguridad quiere frenar
todo para revisar. Hoy la mayoría aplica seguridad recién en la \textbf{última
etapa} (operación). DevSecOps propone el \textbf{Shift Left}: meter la seguridad
lo más a la izquierda posible (planificación/desarrollo).

\prof{Cuanto más temprano detectamos un problema, menos difícil es resolverlo.}

\paragraph{Seguridad por etapa (Shift Left):}
\begin{itemize}
  \item \textbf{Plan \& Develop}: \textbf{Threat modeling} (evaluar amenazas de
    una solución nueva antes de desarrollarla), plugins de seguridad en el IDE,
    \textbf{pre-commit hooks} (que no se expongan secretos), buenas prácticas,
    revisión de PRs.
  \item \textbf{Commit the code / PR}: \textbf{SAST} (análisis estático, p.ej.\
    SonarQube), unit/functional tests de seguridad (foco en \textbf{CWE}),
    control de vulnerabilidades en dependencias, \textbf{pipelines seguros}
    (protección de secretos).
  \item \textbf{Build \& Test}: \textbf{DAST} (análisis dinámico en sandbox),
    validación de configuración cloud / IaC, infrastructure scanning, security
    acceptance testing.
  \item \textbf{Go to Production}: evaluación de la app como un todo incluyendo
    el ambiente de despliegue (Firewalls, WAF, IPS), controles sobre artefactos
    y secretos, \textbf{pentests externos}, \textbf{Red/Blue team}.
  \item \textbf{Operate}: monitoreo continuo de métricas y logs, integración con
    Threat Intelligence, blind pentest, simulación de incidentes, respuesta a
    incidentes, \emph{blameless postmortems}.
\end{itemize}

\paragraph{Testeos de seguridad en el código.}
\begin{itemize}
  \item \textbf{SAST} (Static): analiza el código \emph{sin ejecutarlo}; marca
    malas prácticas, code smells y problemas de seguridad.
  \item \textbf{DAST} (Dynamic): evalúa la app \textbf{en ejecución},
    interactuando como un usuario externo a través de las interfaces estándar
    (APIs, endpoints); busca fallas explotables en red (validación de entrada,
    autenticación). Automatiza pruebas repetibles (escanear puertos, inputs
    inválidos, etc.).
  \item \textbf{SCA} (Software Composition Analysis): analiza
    \textbf{dependencias y librerías} para identificar vulnerabilidades
    (incluso \textbf{transitivas}), problemas de licencia y otros riesgos. Debe
    extenderse a servicios SaaS consumidos.
\end{itemize}

\begin{destacado}
Caso real mencionado: \textbf{Log4Shell} (Log4j, 2021). Una vulnerabilidad
crítica de \emph{ejecución remota de código} en una librería de logging de Java
usadísima. El problema: muchas empresas \textbf{no sabían si la usaban} (era una
dependencia tan estándar que ``te olvidás que está''). De ahí el resurgir del
\textbf{SBOM}.
\end{destacado}

\begin{definicion}{SBOM (Software Bill of Materials)}
Inventario \textbf{formal y legible por máquina} de los componentes de software,
dependencias, información sobre ellos y sus relaciones. Debe ser lo más completo
posible (e indicar dónde no se pueden articular relaciones) e incluye software
open-source y comercial. Permite saber qué componentes usa cada aplicación ante
un incidente.
\end{definicion}

% =====================================================================
\subsection{Protección de la nube}

La configuración de la nube es \textbf{responsabilidad del cliente}. La
\textbf{Infraestructura como Código (IaC)}, p.ej.\ Terraform, permite definir la
configuración de forma reproducible; existen análisis estáticos que detectan
configuraciones inseguras o ineficientes (puertos expuestos, redes abiertas).

Herramientas (productos caros, pensados para empresas con miles de servidores):
\begin{itemize}
  \item \textbf{CSPM (Cloud Security Posture Management)}: busca problemas de
    configuración y permisos; analiza IaC \emph{antes} de desplegar y revisa la
    config actual. Enfoque más \textbf{estático}.
  \item \textbf{CWPP (Cloud Workload Protection Platform)}: protege las
    aplicaciones corriendo en la nube (VMs, contenedores, serverless); analiza
    \textbf{comportamiento}, microsegmentación a nivel de contenedores, threat
    intelligence. Complementa a CSPM.
  \item \textbf{CIEM (Cloud Infrastructure Entitlement Management)}: foco en los
    \textbf{permisos} de los sujetos sobre recursos cloud; busca permisos
    excesivos comparando uso real vs.\ permisos otorgados; visualiza permisos
    entre cuentas.
  \item \textbf{CNAPP (Cloud Native Application Protection Platform)}:
    plataforma que integra lo anterior (visión seguridad proactiva +
    reactiva).
\end{itemize}

% =====================================================================
\subsection{Evaluaciones de seguridad}

\prof{Tomamos todos los recaudos, todas las mejores prácticas\ldots\ y en algún
momento ese sistema va a fallar. No pregunten cuándo, no pregunten cómo, pero va
a fallar.}

La seguridad nunca está asegurada $\Rightarrow$ hay que evaluar
\textbf{constantemente}. Tipos de evaluación:
\begin{itemize}
  \item \textbf{Evaluación de vulnerabilidades}: herramientas
    \textbf{automatizadas} que escanean en busca de vulnerabilidades
    \textbf{conocidas}.
  \item \textbf{Pentesting}: más exhaustivo, lo hacen \textbf{expertos} que
    intentan explotar vulnerabilidades de forma controlada.
  \item \textbf{Evaluación de riesgos}: identificar/evaluar/priorizar riesgos e
    impacto (la seguridad como una amenaza más).
  \item \textbf{Auditoría de seguridad}: revisión sistemática de políticas,
    procedimientos y controles (interna o externa).
  \item \textbf{Evaluación de cumplimiento}: obligatoria para regulaciones
    (\textbf{GDPR, HIPAA, PCI-DSS}, protección de datos personales).
  \item \textbf{Seguridad física}: control de acceso a instalaciones, cámaras,
    protección de archivos físicos y equipos.
  \item \textbf{Seguridad de aplicaciones}: revisión de código, pruebas de
    seguridad y análisis de arquitectura.
\end{itemize}

\subsubsection{Penetration Test (PenTest)}
Objetivo: agarrar el sistema \textbf{funcionando} y probarlo integralmente para
identificar las vulnerabilidades que lo hacen pasar de un estado \textbf{seguro}
a uno \textbf{inseguro}. Se hace con \textbf{alcance y tiempo definidos} (no se
puede probar eternamente).

\paragraph{Tipos según visibilidad del evaluador:}
\begin{itemize}
  \item \textbf{Black Box / Blind}: el evaluador no tiene información previa
    (quizás solo el dominio/IP). Simula un atacante externo; favorece pensar
    ``fuera de la caja'' y descubrir vectores que el sesgo ocultaría. Razón
    práctica: no querer compartir info interna (motivos legales).
  \item \textbf{White Box}: acceso completo (código fuente, diagramas,
    credenciales). Evaluación exhaustiva de vulnerabilidades internas.
  \item \textbf{Gray Box}: intermedio (algún conocimiento, no todo). Simula un
    usuario con acceso limitado (empleado con privilegios restringidos).
  \item \textbf{Double Blind}: Black Box donde el \textbf{atacado tampoco sabe}
    del ejercicio; evalúa además las capacidades de monitoreo y respuesta. Es lo
    más parecido a la realidad (relacionado con Chaos Monkey).
\end{itemize}

\paragraph{Tipos según sistema evaluado:} Red externa (servicios expuestos a
Internet), Red interna (atacante interno: empleado descontento o malware en una
workstation), Redes inalámbricas (NFC, Bluetooth, Wireless), Aplicaciones web
(SQLi, XSS, etc.) e \textbf{Ingeniería social} (phishing, smishing).

\paragraph{Etapas (6 pasos):} 1) Preparación (objetivo, alcance, tiempo, nivel
de riesgo aceptado) $\to$ 2) Armado del plan de ataque (herramientas, ubicación)
$\to$ 3) Identificar el equipo $\to$ 4) Definir el objetivo a alcanzar $\to$
5) Ejecutar el ataque (termina al lograr el objetivo o al vencerse el tiempo)
$\to$ 6) Integrar los hallazgos y reiniciar el ciclo (pivotear/escalar con lo
obtenido).

\subsubsection{Red / Blue / Purple Team}
\begin{itemize}
  \item \textbf{Red Team} (ofensivo): piensa y actúa como atacante
    malintencionado de forma controlada y ética --- pentesting, ingeniería
    social, assessment de vulnerabilidades.
  \item \textbf{Blue Team} (defensivo): implementa y mantiene controles
    (firewalls, IDS, gestión de parches y configuraciones), monitorea y responde
    a incidentes.
  \item \textbf{Purple Team}: enfoque colaborativo que combina ambos para
    mejorar la postura de seguridad (fomenta comunicación en vez de aislamiento).
\end{itemize}

\subsubsection{Breach and Attack Simulation (BAS)}
Categoría de herramientas/técnicas que evalúan la efectividad de las defensas
simulando las \textbf{tácticas, técnicas y procedimientos (TTP)} de un atacante
real. Objetivo: identificar vulnerabilidades \textbf{y} medir la capacidad de la
organización para \textbf{detectar y reaccionar}. Se ejecutan de forma
\textbf{recurrente} en distintos puntos. Relacionado: \emph{Breach Attack
Simulation}/juegos de guerra (simular un escenario desfavorable puntual ---
``¿qué pasa si un atacante obtiene una ruta de admin del cloud?'' --- y analizar
qué tácticas de defensa aplicar); es barato (juntar gente y plantear el
escenario, a veces en ambiente controlado).

\end{document}
